\documentclass[10pt]{beamer-control}
\usepackage{beamer-control-singlefile}
\INCLUDEONLY{The Nyquist Criterion}
\begin{document}
\CONCEPT{The Nyquist Criterion}

\begin{SUMMARY}
\begin{itemize}
\item Nyquist plot
\item Nyquist criterion
\end{itemize}
\vfill References:
\begin{itemize}
\item \astrom{§10.2}
\end{itemize}
\end{SUMMARY}


\SUBCONCEPT{Review of stability}

\begin{frame}{Definitions}
\begin{description}
  \item[Plant] $P$
  \item[Controller] $C$
  \item[Loop transfer function] $CP$

  \item[Open loop stability]\

  When the poles of $CP$ are in the LHS of the $s$-plane.

  \item[Closed loop stability]\

  When the poles of $\frac{CP}{1+CP}$ are in the LHS of the $s$-plane.

\end{description}
\end{frame}


\begin{frame}{Examples of stability with gain}
Consider the simple loop transfer function:
\begin{align*}
CP &= \frac{k(s-2)}{(s+4)^2} & \frac{CP}{1+CP} &= T = \frac{k(s-2)}{(s+4)^2+k(s-2)}
\end{align*}
Always open loop stable.
\begin{align*}
\text{For } k= 1:\qquad T &= \frac{(s-2)}{s^2+9s+14} \qquad \text{stable poles}\\
\text{For } k=10:\qquad T &= \frac{10(s-2)}{s^2+18s-4} \qquad \text{\emph{unstable} poles}
\end{align*}
Here, we could analytically determine $k=8$ is the threshold between stability and instability, but not (easily) for more complex systems.
\end{frame}


\SUBCONCEPT{Nyquist plot}

\begin{frame}{Frequency response graphically}
\begin{itemize}
\item The dynamics of a linear system may be represented by its frequency response and graphically illustrated by a Bode plot
\item We may also graphically represent the frequency response on one graph by using a Nyquist plot, plotted on the complex plane
\item The Nyquist plot of a loop transfer function $L(s)$ is formed by tracing $s\in\mathbb{C}$ around the Nyquist contour $\Gamma$ (given by the imaginary axis combined with an arc at infinity connecting the endpoints of the imaginary axis)
\item At any point on the Nyquist plot corresponding to a frequency $s=\jj\omega$, the distance away from the origin gives the magnitude of the loop transfer function $L(s)$ and the angle with the positive real axis gives the phase
\end{itemize}
\end{frame}


\begin{frame}
\frametitle{Nyquist contour to the Nyquist plot (cf.\ Fig.\,10.4)}

\end{frame}

\begin{frame}{Nyquist plot}
\begin{figure}
	\centering
	\includegraphics[width=0.9\linewidth]{figure10.5}
	\\
	\textbf{Figure 10.5:} Nyquist plot for $L(s) = {1}/{(s+0.6)^3}$.
\end{figure}
\end{frame}


\SUBCONCEPT{Nyquist criterion}

\begin{frame}{Condition for stability - Nyquist criterion}
\begin{itemize}
\item Let $L(s)$ be the loop transfer function for a negative feedback system and assume (initially) that $L$ has no poles in the closed right half-plane except possibly at the origin
\item Then the closed loop system 
\[
G_{cl}=\frac{L(s)}{1+L(s)}
\]
is stable iff the image of $L$ along the closed contour $\Gamma$ has no net encirclements of the critical point $-1$ in the complex plane
\item The number of encirclements is called the winding number
\item This criterion tells us how the stability of a system is influenced by changing the controller parameters
\end{itemize}
\end{frame}

\begin{frame}{Nyquist plot --- simple stability example}
	\centering
	\includegraphics[width=0.6\linewidth]{figure10.5}
\end{frame}

\begin{frame}{Example}
\begin{itemize}
	\item Consider the transfer function 
	\[L(s) = \frac{k}{s(s+1)^2}\] 
	with nominal gain value $k=1$
	\item The system has a single pole at $s=0$ and a double pole at $s=-1$
	\item The Bode plot thus has the slope $-1$ for low frequencies, and at the double pole the slope changes to $-3$
	\item The phase curve starts at $-90^\circ$ for low frequencies, is \ang{-180} at $\omega=1$ and is \ang{-270} at high frequencies
\end{itemize}
\end{frame}


\begin{frame}{Example continued}
	\begin{itemize}
		\item The Nyquist curve intersects the negative real axis at $\omega=1$ with $L(\ii)=-0.5$ for $k=1$
		\item The Nyquist curve only encircles the critical point $s=-1$ when the gain is $k\geq 2$ as $L(\ii)=-\tfrac{k}{2}$
	\end{itemize}
\begin{figure}
	\centering
	\includegraphics[width=0.9\linewidth]{figure10.6}
	\\
	\textbf{Figure 10.6:} Sketching Nyquist and Bode plots.
\end{figure}
\end{frame}


\SUBCONCEPT{Generalised Nyquist criterion}

\begin{frame}{Contour in PZMAP = different contour in freq resp}
  \includegraphics[width=\linewidth]{encirc1}
\end{frame}
\begin{frame}{Contour in PZMAP = different contour in freq resp}
  \includegraphics[width=\linewidth]{encirc3}
\end{frame}
\begin{frame}{Contour around pole}
  \includegraphics[width=\linewidth]{encirc2}
\end{frame}
\begin{frame}{Contour around zero}
  \includegraphics[width=\linewidth]{encirc4}
\end{frame}

\begin{frame}{The Cauchy criterion}
When taking a closed contour in the complex plane, and mapping it through a complex function $L(s)$:
\[
  n_N = n_Z - n_P
\]
where:
\begin{itemize}
\item $n_N =$ number of times the mapping of $L(s)$ encircles the origin
\item $n_Z =$ number of zeros of $L(s)$ enclosed by contour
\item $n_P =$ number of poles of $L(s)$ enclosed by contour
\end{itemize}
Note that positive $n_N$ = same direction as the original contour
\end{frame}

\begin{frame}{Poles/zeros to encirclement and vice versa}
  \begin{itemize}
    \item  These examples show if we know the poles and zeros in a contour we can determine the number of encirclements
    \item  The opposite also holds!
    \item  By looking at encirclements we can predict the poles and zeros found in the region
    \item  Here's the trick:
    \begin{itemize}
    \item  Our contour is defined as the entire RHS of the $s$-plane
    \item  Our mapping is $s=\jj\ww$ for $\omega\in(-\infty,+\infty)$
    \item  And one more important thing\dots
    \end{itemize}
  \end{itemize}
\end{frame}

\begin{frame}{Contour along imaginary axis}
  \includegraphics[width=\linewidth]{encirc5}
\end{frame}
\begin{frame}{As it grows larger\dots}
  \includegraphics<1>[width=\linewidth]{encirc6}
  \includegraphics<2>[width=\linewidth]{encirc7}
  \includegraphics<3>[width=\linewidth]{encirc8}
\end{frame}

\begin{frame}{That final important thing}{The critical point $(-1,0)$}
  \begin{itemize}
    \item  $L$ encirclement of origin: $L$ has RHS poles and zeros
    \begin{itemize}
      \item  Don't really care --- we usually know when $L$ is unstable!
    \end{itemize}
    \bigskip
    \item  $L$ encirclement of $(-1,0)$ $\equiv$ $(1+L)$ encirclement of origin
    \bigskip
    \item  If $1+L$ has RHS zeros, $\frac{L}{1+L}$ will have RHS poles
    \begin{itemize}
    \item  This is what we want to know!
    \item  Use encirclement to find RHS zeros of $1+L$
    \item  Poles of $1+L$ are same as the poles of $L$ (\emph{think about it})
    \end{itemize}
  \end{itemize}
\end{frame}

\begin{frame}{How many unstable CL poles?}
Cauchy Criterion reformulated as the Nyquist Stability Criterion:
\FORMULA{n_Z = n_P + n_N}
\begin{equation}
\theFORMULA
\end{equation}
For loop transfer function $L$:
\begin{itemize}
\item $n_Z=$ number of RHS CL poles (of $\frac{L}{1+L}$) \hfill \makebox[0.4\linewidth][r]{\small\em --- or \textbf{z}eros of $1+L$}\smallskip

\item $n_P=$ number of RHS OL poles (of $L$) \hfill \makebox[0.4\linewidth][r]{\small\em --- know this from the TF}\smallskip

\item $n_N=$ number of CW encirc.\ of $(-1,0)$ \hfill \makebox[0.4\linewidth][r]{\small\em --- get this from Nyquist}
\end{itemize}

\end{frame}

\begin{frame}{Non-negative number of poles and zeros}
$n_Z\ge 0$ and $n_P\ge 0$ imply the following scenarios:
\begin{multicols}{2}
$n_P=0$:
\begin{align*}
n_N=0 &\implies n_Z=0 \\
n_N>0 &\implies n_Z>0
\end{align*}
Avoid encirclement for stability.
\bigskip

$n_P>0$:
\begin{align*}
n_N=-n_P &\implies n_Z=0 \\
n_N\ge 0 &\implies n_Z>0 \\
0>n_N>-n_P &\implies n_Z>0
\end{align*}
Need $n_P$ CCW encirclements for stability.
\end{multicols}

Finally, note that \emph{any} CW encirclements always means CL instability.
\end{frame}



\SUMMARYFRAME
\FINALE

\end{document}
